\textbf{ABAQ}\\
Plataforma de Servicio Social

\begin{center}\rule{0.5\linewidth}{0.5pt}\end{center}

\textbf{Reporte de Avance No.~1}\\
Presentación inicial del prototipo y definición de requerimientos

\begin{center}\rule{0.5\linewidth}{0.5pt}\end{center}

{\def\LTcaptype{none} % do not increment counter
\begin{longtable}[]{@{}rl@{}}
\toprule\noalign{}
\endhead
\bottomrule\noalign{}
\endlastfoot
\textbf{Fecha:} & 19 de agosto de 2026 \\
\textbf{Versión:} & v1.0.0 \\
\textbf{Estado:} & En planeación / Prototipado \\
\textbf{Preparado por:} & Equipo de Desarrollo ABAQ \\
\end{longtable}
}

Documento confidencial --- Solo para uso interno del cliente

\section*{Resumen ejecutivo}\label{resumen-ejecutivo}
\addcontentsline{toc}{section}{Resumen ejecutivo}

Este documento presenta el primer reporte de avance del proyecto ABAQ.
Durante este período se llevó a cabo la presentación oficial del primer
prototipo funcional de la plataforma a la asociación ABAQ Querétaro,
logrando alinear las expectativas sobre el diseño y la arquitectura del
sistema. Se acordó el desarrollo de un sistema con accesos separados
para administradores y estudiantes, garantizando la seguridad de la
información. Finalmente, se negoció un esquema de trabajo de 8 entregas
quincenales para la liberación del servicio social.

\section{Introducción}\label{introducciuxf3n}

Este documento resume los acuerdos y requerimientos recabados en la
primera reunión oficial de seguimiento (19 de agosto de 2026) entre el
equipo de desarrollo de software y ABAQ Querétaro. El propósito
principal de la sesión fue validar el rumbo técnico del proyecto y
formalizar las condiciones operativas.

\section{Cambios en el Servidor
(Backend)}\label{cambios-en-el-servidor-backend}

\subsection{Resumen general}\label{resumen-general}

En esta etapa, se definió la arquitectura base para la seguridad y
separación de perfiles en la plataforma.

\subsection{Protección de rutas con
autenticación}\label{protecciuxf3n-de-rutas-con-autenticaciuxf3n}

Se acordó que el sistema requerirá autenticación obligatoria. El acceso
de los estudiantes a módulos administrativos estará bloqueado a nivel de
sistema.

\subsection{Mejoras en la
configuración}\label{mejoras-en-la-configuraciuxf3n}

\emph{Pendiente por implementar / Falta de información para esta etapa.}

\begin{longtable}[]{@{}ll@{}}
\caption{Resumen de cambios en el servidor
(Backend)}\label{tab:backend-changes}\tabularnewline
\toprule\noalign{}
\textbf{Área} & Descripción del cambio \\
\midrule\noalign{}
\endfirsthead
\toprule\noalign{}
\textbf{Área} & Descripción del cambio \\
\midrule\noalign{}
\endhead
\bottomrule\noalign{}
\endlastfoot
\textbf{Seguridad} & Definición de separación estricta entre perfiles de
Administrador y Estudiante. \\
\textbf{Infraestructura} & Validación del almacenamiento en MongoDB (500
MB) como escalable y suficiente. \\
\end{longtable}

\section{Cambios en la Aplicación Web
(Frontend)}\label{cambios-en-la-aplicaciuxf3n-web-frontend}

\subsection{Resumen general}\label{resumen-general-1}

Se presentó el prototipo visual y se recogieron los requerimientos para
los siguientes módulos del sistema.

\subsection{Panel de Administración}\label{panel-de-administraciuxf3n}

\emph{Pendiente. Se presentará la distinción de vistas en la próxima
sesión.}

\subsection{Vista de Detalle de
Estudiante}\label{vista-de-detalle-de-estudiante}

\emph{Falta de información / Pendiente de implementación.}

\subsection{Vista de Lista de
Estudiantes}\label{vista-de-lista-de-estudiantes}

\emph{Pendiente. Se acordó alimentar la plataforma con datos de prueba
realistas para la próxima sesión.}

\subsection{Lista de Actividades}\label{lista-de-actividades}

Se establecieron los requerimientos para la gestión de campañas (ej.
jornadas de esterilización) y el registro de asistencia con control de
entradas y salidas.

\subsection{Panel del Estudiante (Vista de
inicio)}\label{panel-del-estudiante-vista-de-inicio}

\emph{Pendiente por presentar en la siguiente revisión.}

\subsection{Carga de Archivos}\label{carga-de-archivos}

Se solicitó la capacidad de automatizar la generación de constancias de
término o cartas de participación desde el sistema de forma digital.

\begin{longtable}[]{@{}ll@{}}
\caption{Resumen de cambios en la aplicación web
(Frontend)}\label{tab:frontend-changes}\tabularnewline
\toprule\noalign{}
\textbf{Área} & Descripción del cambio \\
\midrule\noalign{}
\endfirsthead
\toprule\noalign{}
\textbf{Área} & Descripción del cambio \\
\midrule\noalign{}
\endhead
\bottomrule\noalign{}
\endlastfoot
\textbf{Prototipo} & Presentación y validación del diseño visual inicial
con ABAQ Querétaro. \\
\textbf{Requerimientos} & Definición de funcionalidades para campañas y
generación de constancias. \\
\end{longtable}

\section{Evidencia Visual --- Capturas de
Pantalla}\label{evidencia-visual-capturas-de-pantalla}

\emph{Falta de información. Pendiente de anexar capturas oficiales de
las pantallas mostradas en esta sesión.}

\section{Estado Actual del Proyecto}\label{estado-actual-del-proyecto}

\emph{En Prototipado y Negociación. Se validó el prototipo funcional
inicial y se acordó formalmente la liberación del servicio social a
cambio de la entrega del sistema completo para diciembre de 2026.}

\section{Próximos Pasos}\label{pruxf3ximos-pasos}

\begin{enumerate}
\def\labelenumi{\arabic{enumi}.}
\tightlist
\item
  \textbf{Vistas Separadas:} Desarrollar y presentar el panel dividido
  entre Administrador y Estudiante.
\item
  \textbf{Datos de Prueba:} Alimentar la plataforma con registros de
  prueba realistas.
\item
  \textbf{Documentación Formal:} Elaborar un documento formal que
  detalle la propuesta técnica y el calendario de entregas.
\end{enumerate}

\begin{center}\rule{0.5\linewidth}{0.5pt}\end{center}

\hfill\break
ABAQ --- Reporte de Avance No.~1 · v1.0.0 · Agosto 2026
